% Template LaTeX para Informes de Proyectos de Grado de Computación
% InCo, Facultad de Ingeniería, Universidad de la República
% 06/2023

\documentclass[oneside]{prgrado}



%%%%%%%%%%%%%%%%%%%%%%%%%%%%%%%%%%%%%%%%%%%%%%%%%%%%%%%%%%%%%%%%%%%%%%%%%%%%%%%%%%%%%%%%%%%%%%%%
% Datos del Proyecto:
%%%%%%%%%%%%%%%%%%%%%%%%%%%%%%%%%%%%%%%%%%%%%%%%%%%%%%%%%%%%%%%%%%%%%%%%%%%%%%%%%%%%%%%%%%%%%%%%

% título del proyecto (debe escribirse en minúscula, con excepción de la
% letra inicial de la primera palabra y nombres propios)
\title{Análisis de tópicos en sitios con participación uruguaya}                   

% autores
\author{Juan Manuel Pacheco, Ignacio Perez y Cristhian Rossi}  

% fecha de la defensa
\date{\today}                                 

% supervisor
\supervisor{Juan José Prada}                

% cosupervisor (comentar si no tiene)
\cosupervisor{Santiago Góngora}

% licencia creative commons del documento
% opciones: by, by-sa, by-nd, by-nc, by-nc-sa y by-nc-nd
% opción por defecto: by-nc-nd versión 4.0
\cclicense{by-nc}{4.0}

%%%%%%%%%%%%%%%%%%%%%%%%%%%%%%%%%%%%%%%%%%%%%%%%%%%%%%%%%%%%%%%%%%%%%%%%%%%%%%%%%%%%%%%%%%%%%%%%


%links con colores 
\hypersetup{
  colorlinks   = true,
  urlcolor     = blue,
  linkcolor    = blue,
  citecolor    = red
}


\begin{document}

%%%%%%%%%%%%%%%%%%%%%%%%%%%%%%%%%%%%%%%%%%%%%%%%%%%%%%%%%%%%%%%%%%%%%%%%%%%%%%%%%%%%%%%%%%%%%%%%
% Parte inicial
%%%%%%%%%%%%%%%%%%%%%%%%%%%%%%%%%%%%%%%%%%%%%%%%%%%%%%%%%%%%%%%%%%%%%%%%%%%%%%%%%%%%%%%%%%%%%%%%

\frontmatter % numeración en romanos, capitulos sin numerar

% carátula
\maketitle


%%%%%%%%%%%%%%%%%%%%%%%%%%%%%%%%%%%%%%%%%%%%%%%%%%%%%%%%%%%%%%%%%%%%%%%%%%%%%%%%%%%%%%%%%%%%%%%%

% agradecimientos
\chapter*{Agradecimientos}

Además de mi perro, quiero agradecerle a la facultad y a Claudio. Una por hacerme el trabajo en este artículo más fácil, y la otra por hacérmelo mucho más difícil. 


%%%%%%%%%%%%%%%%%%%%%%%%%%%%%%%%%%%%%%%%%%%%%%%%%%%%%%%%%%%%%%%%%%%%%%%%%%%%%%%%%%%%%%%%%%%%%%%%

% resumen
\chapter*{Resumen}

[Placeholder]

El resumen (200-500 palabras) debe dar una idea completa de todo el
proyecto, mencionando claramente los formalismos, técnicas, herramientas y lenguajes
utilizados. No debe limitarse a describir el problema abordado, sino que debe describir
la solución del problema, con una evaluación de la misma. No debe incluir referencias
bibliográficas ni referencias a otras partes del informe. Tampoco debe utilizar
acrónimos sin explicar su significado.

[Que poner]

En la actualidad, cada persona tiene la oportunidad de expresar su opinión en un mundo más globalizado que nunca; cada día se generan en promedio XXX mensajes en redes sociales. Por ello, desarrollar herramientas capaces de procesar todas estas opiniones constituye un objetivo más justificado que nunca. Limitar el estudio al entorno uruguayo nos brinda la oportunidad de enfocarnos y ejercer un mayor control sobre los temas discutidos en una región tan particular como la República Oriental del Uruguay, además de permitirnos explorar el procesamiento de un lenguaje más ruidoso y abierto a los uruguayismos propios de la zona.

Exploramos y experimentamos con diversos métodos de preprocesamiento con el propósito de extraer la mayor utilidad de cada mensaje, así como con diferentes modelos para determinar cuál se adapta mejor a la problemática de trabajar con documentos de un tamaño tan reducido como el de los mensajes en redes sociales.

\hfill \break
\keywords{Procesamiento del Lenguaje Natural, Aprendizaje Automático, Modelado de Tópicos, Redes Sociales}

%%%%%%%%%%%%%%%%%%%%%%%%%%%%%%%%%%%%%%%%%%%%%%%%%%%%%%%%%%%%%%%%%%%%%%%%%%%%%%%%%%%%%%%%%%%%%%%%

% índice
\tableofcontents
\newpage


%%%%%%%%%%%%%%%%%%%%%%%%%%%%%%%%%%%%%%%%%%%%%%%%%%%%%%%%%%%%%%%%%%%%%%%%%%%%%%%%%%%%%%%%%%%%%%%%
% Parte central
%%%%%%%%%%%%%%%%%%%%%%%%%%%%%%%%%%%%%%%%%%%%%%%%%%%%%%%%%%%%%%%%%%%%%%%%%%%%%%%%%%%%%%%%%%%%%%%%

\mainmatter %numeración en arábicos, numerar capítulos 

% introducción
\chapter{Introducción}

[Placeholder]

Aquí se motiva el trabajo, se plantea y define el problema, se deja claro cuales son los objetivos (general del proyecto, si correspondiese o si está inmerso en un proyecto de mayor alcance, y los específicos), se plantean los resultados esperados, se establecen resumidamente las conclusiones y se describe la
organización general del documento.

Este y los siguientes capítulos pueden incluir referencias, algunos ejemplos de referencias pueden ser \cite{CitekeyArticle}, \cite{CitekeyBook}, \cite{CitekeyInproceedings}, \cite{CitekeyManual} y \cite{CitekeyMisc}. Se sugiere citar usando el formato APA.

[Que poner]

Esta introccion para mi amerita incluir aspectos tales como el enfoque principal de investigacion del proyecto tales como es el analisis de textos usando PLN y las dificultades que dispone de por si y las adicionales resultantes de trabajar con mensajes del entorno de redes sociales.
Enfoques por lo que va a ir esta tesis, trabajar con diferentes tipos de prepocesamientos (conbinaciones de los mismos), diferentes modelos de procesamientos de textos para devolver que temas fueron devueltos.

Y establecer milestones sobre que es posible utilizar en el contexto de estudio de mensajes en contextos en redes sociales en la region rioplatense.

\section{Problema y motivación}

Poner algo sobre el flujo de mensajes y como es demasiado para que una persona lo vea, impulsar la necesidad de tener herramientas de estas caracteristicas, 

\section{Objetivos}

Investigar como hacer con el lenguaje ruidoso de redes sociales y como tratarlo, que modelos los pueden procesar mejor estos textos y como puedo efectivamente tratar de procesar el flujo de los mismo y tener una idea aproximada de lo hablado 


%%%%%%%%%%%%%%%%%%%%%%%%%%%%%%%%%%%%%%%%%%%%%%%%%%%%%%%%%%%%%%%%%%%%%%%%%%%%%%%%%%%%%%%%%%%%%%%%

% antecedentes
\chapter{Revision de antecedentes}

[Placeholder]

Revisión de antecedentes (ya sea productos, procesos, publicaciones, etc., a nivel académico o comercial) en el tema del trabajo. Puede incluir además (si es necesario) una breve introducción a los conceptos necesarios para entender el trabajo.

[Que poner]

En este capítulo se explican algunos conceptos, técnicas y tareas vinculadas a la Inteligencia Artificial, al Procesamiento de Lenguaje Natural y otros conceptos necesarios para enmarcar el proyecto. El propósito de este capítulo es contextualizar al lector en la teoría referenciada en las secciones posteriores de este informe.

En la sección ~\ref{sec:CAP2_IA} se introduce el concepto de Inteligencia Artificial. Luego, en la sección ~\ref{sec:CAP2_AA} se describen los conceptos principales asociados con el Aprendizaje Automático, donde se explican las principales etapas y técnicas utilizadas en este proyecto. En la sección ~\ref{sec:CAP2_PLN} se presenta la disciplina del Procesamiento de Lenguaje Natural, presentando los modelos en los que se basa la solución de este proyecto. Por último, en la sección ~\ref{sec:CAP2_Extra} se abordan otros conceptos útiles y complementarios para el proyecto, como las Bases de Datos y las Interfaces Web.

El lector que ya cuente con conocimientos sobre estos temas puede optar por saltear la lectura de este capítulo.

\section{Inteligencia Artificial}\label{sec:CAP2_IA}

\section{Aprendizaje Automático}\label{sec:CAP2_AA}

\subsection{Aprendizaje Supervisado}

\subsection{Aprendizaje No Supervisado}

\subsection{Entrenamiento, validación y verificación}

\subsection{Medidas de evaluación}
ACA o mas a bajo?
Porque si usamos Matrices de confucion capaz la idea seria ponerlo aca, pero si va a ser el metodo casero hasta mas abajo prodria ir porque es custom para PLN

\section{Procesamiento de Lenguaje Natural}\label{sec:CAP2_PLN}

\subsection{Preprocesamiento}

\subsubsection{Word embeddings}

\subsection{Modelos Utilizados}

\subsubsection{Latent Dirichlet Allocation}

\subsubsection{Bert Topic}

\subsubsection{Biterm Topic Modeling}

\subsubsection{Biterm Dirichlet Process}

\subsection{Medidas de evaluación}
O aca qsy

\section{Conceptos Complementarios}\label{sec:CAP2_Extra}

\subsection{Base de Datos Relacional}

\subsection{Contenedor/Docker}

\subsection{Web API}

%%%%%%%%%%%%%%%%%%%%%%%%%%%%%%%%%%%%%%%%%%%%%%%%%%%%%%%%%%%%%%%%%%%%%%%%%%%%%%%%%%%%%%%%%%%%%%%%

% parte central
\chapter{Parte Central}
Name: Identificación dinamica de tópicos
[Placeholder]

La parte central del trabajo refiere a lo que es producción propia o aporte del
proyecto de grado, incluyendo las decisiones tomadas. Por ejemplo, puede incluir
los requerimientos, el análisis y el diseño de la solución. Si el proyecto tiene una
implementación, debe describirse en términos de decisiones tomadas en ese sentido.
Los detalles de programación se dejan para los anexos.

Se pueden incluir figuras y tablas en el documento, las mismas deben estar referenciadas en el texto. Por ejemplo, la Figura muestra los logos de Facultad de Ingeniería y de la Universidad de la República, mientras que la Tabla tiene números aleatorios.

[Que poner]

\section{Serie de pasos}

\subsection{Obtención de mensajes}

\subsection{Preprocesamiento de mensajes}

\subsection{Procesamiento de mensajes}

\subsection{Evaluación}

\subsection{Visualizacion de Topicos}


\section{Aplicación}

\subsection{Orquestrador}

\subsection{Usuario}


%%%%%%%%%%%%%%%%%%%%%%%%%%%%%%%%%%%%%%%%%%%%%%%%%%%%%%%%%%%%%%%%%%%%%%%%%%%%%%%%%%%%%%%%%%%%%%%%

% experimentación
\chapter{Experimentación}
Puede ser necesario incluir un capítulo de Experimentación, incluyendo las pruebas realizadas (casos de prueba) y los resultados obtenidos con su respectivo análisis, que puede incluir comparaciones. 

%%%%%%%%%%%%%%%%%%%%%%%%%%%%%%%%%%%%%%%%%%%%%%%%%%%%%%%%%%%%%%%%%%%%%%%%%%%%%%%%%%%%%%%%%%%%%%%%

%conclusiones
\chapter{Conclusiones y Trabajo Futuro}

En este capítulo se evalúan los resultados alcanzados y
dificultades encontradas, se establece lo que se planteó hacer y lo que se hizo
realmente, cuales fueron los aportes, se muestran posibles extensiones al trabajo, se
realiza una autocrítica de lo que se hizo y lo que faltó (por problemas de tiempo,
recursos, cómo se puede continuar, qué cosas hacer, prioridades, etc.) y se incluye
información sobre la gestión del proyecto, si aplica

%%%%%%%%%%%%%%%%%%%%%%%%%%%%%%%%%%%%%%%%%%%%%%%%%%%%%%%%%%%%%%%%%%%%%%%%%%%%%%%%%%%%%%%%%%%%%%%%
% Parte final
%%%%%%%%%%%%%%%%%%%%%%%%%%%%%%%%%%%%%%%%%%%%%%%%%%%%%%%%%%%%%%%%%%%%%%%%%%%%%%%%%%%%%%%%%%%%%%%%

{ % inicio backmatter

\backmatter % no numerar capítulos 

% referencias bibliográficas

\newpage
\bibliographystyle{apacite}
\bibliography{referencias}

} % fin backmatter

%%%%%%%%%%%%%%%%%%%%%%%%%%%%%%%%%%%%%%%%%%%%%%%%%%%%%%%%%%%%%%%%%%%%%%%%%%%%%%%%%%%%%%%%%%%%%%%%

% anexos
\begin{appendix}

\chapter{Anexo 1}

Los anexos contienen información adjunta al proyecto pero que no es fundamental
para entender el trabajo. Por ejemplo, determinado material de los antecedentes o la
implementación, que en el cuerpo principal del informe se encuentre resumido, aquí
puede presentarse de forma completa. En caso de proyectos de desarrollo de
software, se debería incluir un manual de usuario.

\section{Sección del Anexo}

Los anexos pueden tener secciones.

\end{appendix}

%%%%%%%%%%%%%%%%%%%%%%%%%%%%%%%%%%%%%%%%%%%%%%%%%%%%%%%%%%%%%%%%%%%%%%%%%%%%%%%%%%%%%%%%%%%%%%%%

\end{document}
